% 07_limitations.tex: Section 6.
\section{Limitations}\label{sec:limitations}

The results in this report support a limited conclusion, and I list here the reasons for this limit. I order them by how strongly they affect the interpretation of the results.

\paragraph{One test month.} All test results come from December 2025, which contains 967 flights. The bootstrap intervals in \cref{tab:metrics} show that the metrics vary with the choice of test days, and a different month would give different values. The test month is also not typical: its delayed share of 54.1\% is much higher than the 40.9\% of the training months and the 35.5\% of the validation month. Both models flagged fewer delays than occurred (\cref{sec:confusion}), and the accuracy of the majority baseline depends on the delayed share as well. I therefore treat the absolute values of accuracy, precision, and recall as specific to December, and I give more weight to the ROC AUC, which does not depend on the delayed share.

\paragraph{Fixed threshold.} I used the threshold of 0.5 for all reported results (\cref{sec:rule}). The threshold analysis shows that another threshold would give a different balance of false alarms and missed delays (\cref{sec:threshold}). I did not choose a threshold from the costs of the two errors, because I do not have this information, and a threshold chosen on the test month would not be a fair estimate of the performance on new data.

\paragraph{Schedule-only features.} The models use seven variables that are part of the schedule (\cref{sec:features}). I excluded all information that becomes known at or after arrival, which is necessary for a valid prediction, but it also means that the models cannot see the causes of delays that an operator can observe, such as weather, the delay of the inbound flight, or the situation on the apron. The moderate ROC AUC of 0.74 to 0.75 is therefore not surprising. It shows how much the schedule alone says about the delay, and it is not an upper limit for what a better-informed system could reach.

\paragraph{Rescheduled flights.} The ledger contains 118 flights whose remark marks them as rescheduled, and for these flights, STA may not describe the plan that the airline actually followed. The target is the difference to the recorded STA, so the delay of such a flight may be overstated or understated, and I cannot verify this from the data. I kept these flights in the main analysis and tested their removal in \cref{sec:robustness}. The results changed little, but the test only covers the flights that carry the remark, and other changes of the plan may not be marked.

\paragraph{Scope of the uncertainty estimates.} The bootstrap resamples the 31 test days, and the models are not refitted (\cref{sec:bootstrap}). The intervals therefore measure how the metrics depend on the test days. They do not measure the uncertainty from the training data, from the selection of the settings on a single validation month of 902 flights, or from the seed. The comparison of the two models is also limited: the settings of both were selected within a range of about 0.01 macro F1 on the validation month, and the paired difference on the test month is small. I do not claim that one model is better in general. In addition, the bootstrap, the threshold analysis, and the robustness analysis reuse the December data as a diagnosis, although none of them changed a setting.

\paragraph{Small groups.} The error analysis in \cref{sec:errors} rests on small groups. Most airlines have fewer than 100 test flights, and the pooled group of smaller airlines mixes operators with different behavior. I did not compute intervals for these groups, and their error rates should be read as descriptions of the test month and not as properties of the airlines. The counts of false positives and false negatives are derived from the saved tables and not taken from the notebook directly.

\paragraph{Probabilities.} Both models use balanced class weights, so their scores are not calibrated probabilities (\cref{sec:models}). I used them only to rank flights and to apply a threshold. I did not evaluate calibration, and the scores should not be interpreted as the chance of a delay.

\paragraph{Non-public data.} The ledger is restricted and cannot be shared (\cref{sec:data}). Other researchers can inspect the code and the aggregated results in the repository, but they cannot rerun the analysis without the ledger, which is held by AIBD LAS and would need to be requested from AIBD LAS. Airline names are anonymized in the saved tables, so the airline-level results cannot be matched to real operators outside the project.

\paragraph{Generalizability.} The data cover one airport and one calendar year. The models learned the patterns of AIBD in 2025, such as the typical delay of an airline at a given hour, and these patterns can change with a new schedule, new airlines, or new operating conditions. A level of performance like the one in this report should not be expected at another airport or in another year without a new evaluation. Because the training months come before the test month, the setup does reflect the use of a model on future flights at AIBD, but I tested it for one month only.
